% Slides 81--105
\sectionframe{05 / 5 April 2026}{Swaptions become calibrated models, not flat-vol examples}{Commit 31df798; portfolio.py calibrated short-rate model path.}

\chronology{5 Apr 2026}{Curve analytics and calibrated Bermudans arrive}
{Commit \texttt{31df798} adds about 2,082 lines across pricing, model documentation, dashboards, routes, and tests.}
{Hull-White calibration replaces default model parameters.}
{G2++ appears as a comparison model and the QuantLib data-model page exposes construction details.}
{The project begins to separate market fit, model choice, and numerical engine.}
{Commit 31df798.}

\lesson{European swaptions are calibration instruments}
{A diagonal point is identified by option expiry and remaining swap tenor, such as 2Y x 8Y.}
{The market normal volatility is converted into a helper price.}
{The model price is compared with the same helper after assigning a calibration engine.}
{This aligns calibration targets with the future swaps relevant to Bermudan exercise.}
{portfolio.py \_create\_swaption\_helper and diagonal pillars.}

\lesson{Jamshidian decomposition provides a Hull-White benchmark}
{A one-factor coupon-bond option can be decomposed into options on zero-coupon bonds.}
{QuantLib's \texttt{JamshidianSwaptionEngine} prices calibration helpers analytically under Hull-White.}
{Fast helper valuation makes iterative calibration practical.}
{The final Bermudan still uses a tree because multiple exercise dates require backward induction.}
{QuantLib JamshidianSwaptionEngine; portfolio.py.}

\lesson{Mean reversion is fixed; volatility is calibrated}
{The UI supplies Hull-White $a$ as a governed input.}
{The calibration mask fixes $a$ and solves the single constant $\sigma$.}
{Levenberg-Marquardt minimizes the selected helper errors under explicit end criteria.}
{Fixing one parameter avoids trying to identify both $a$ and $\sigma$ from a narrow diagonal.}
{portfolio.py \_build\_hull\_white\_model.}

\begin{frame}{The calibration problem is constrained nonlinear least squares}
  \[
    \widehat\sigma=\arg\min_{\sigma>0}
    \sum_{i=1}^{N} w_i\left(\varepsilon_i(a,\sigma)\right)^2,
    \qquad a\ \text{fixed}.
  \]
  \begin{itemize}\setlength\itemsep{0.8em}
    \item $\varepsilon_i$ may be relative price error or absolute price error.
    \item Initial values, bounds, and termination criteria influence robustness.
    \item Every helper residual is retained for diagnostics and testing.
  \end{itemize}
  \takeaway{Calibration is a reproducible algorithm, not a call to ``fit.''}
  \src{portfolio.py short-rate calibration functions.}
\end{frame}

\lesson{A Bermudan represents a sequence of remaining swaps}
{A 5Y NC2 trade first exercises near year 2 and always ends at the original year-5 maturity.}
{Later decisions enter or cancel progressively shorter remaining swaps.}
{The repository builds labels such as 2Y x 3Y, 3Y x 2Y, and 4Y x 1Y.}
{Holding final maturity fixed is essential; a fixed underlying tenor would describe a different trade.}
{portfolio.py \_bermudan\_schedule\_context; tests/test\_sofr\_curve.py.}

\lesson{Exercise dates are business-day objects}
{The fixed schedule defines underlying start dates and final maturity.}
{Exercise is placed one business day before the corresponding underlying start.}
{QuantLib calendars adjust every date consistently.}
{Tests assert the exact 2028--2031 dates for the benchmark schedule.}
{tests/test\_sofr\_curve.py schedule test.}

\lesson{The diagonal helper strip follows the call schedule}
{Each exercise date maps to the remaining swap tenor.}
{If an expiry is absent from the visible matrix, the repository interpolates between governed neighboring expiries.}
{Source market points are retained and shown in the UI.}
{Calibration lineage connects each model residual to an exercise opportunity.}
{portfolio.py bermudan\_diagonal\_calibration\_pillars.}

\lesson{Tree pricing performs backward induction}
{The short-rate state is discretized on a recombining lattice.}
{At the final exercise date, exercise and continuation values are known at each node.}
{Moving backward, the engine takes the larger holder value at each eligible node.}
{Discounted transition expectations propagate the optimal value to today.}
{QuantLib TreeSwaptionEngine; BermudanSwaption example.}

\begin{frame}{The local exercise rule}
  \[
    V(t_i,x)=\max\left(E(t_i,x),\ C(t_i,x)\right),
    \qquad
    C(t_i,x)=\mathbb E\!\left[D_{i,i+1}V(t_{i+1},X_{i+1})\mid X_i=x\right].
  \]
  \begin{itemize}\setlength\itemsep{0.8em}
    \item The state $x$ is one-dimensional under Hull-White 1F.
    \item Node values share successor states, producing a recombining tree.
    \item Numerical time steps control approximation error between event dates.
  \end{itemize}
  \takeaway{A Bermudan price includes both cash-flow value and the value of a stopping policy.}
  \src{Standard lattice backward induction; QuantLib TreeSwaptionEngine.}
\end{frame}

\lesson{Tree steps are a numerical parameter}
{The repository scales steps with final maturity and enforces a minimum of 50.}
{Too few steps distort exercise and curve fitting on the lattice.}
{More steps increase runtime and may expose date-grid alignment effects.}
{A production model needs an explicit tree-step convergence study.}
{portfolio.py \_bermudan\_tree\_time\_steps.}

\lesson{G2++ adds a second Gaussian rate factor}
{Two mean-reverting factors allow richer yield-curve movements than Hull-White 1F.}
{Volatilities $\sigma$ and $\eta$ are calibrated while mean reversions and factor correlation are fixed in this workbench.}
{QuantLib provides a G2 swaption calibration engine and tree-compatible short-rate model machinery.}
{The comparison reveals model risk without replacing the transparent Hull-White default.}
{portfolio.py \_build\_g2\_model.}

\begin{frame}{G2++ decomposes the short rate into two factors}
  \[
    r_t=\varphi(t)+x_t+y_t,
    \quad dx_t=-a x_tdt+\sigma dW_t^x,
    \quad dy_t=-b y_tdt+\eta dW_t^y,
    \quad dW_t^x dW_t^y=\rho\,dt.
  \]
  \begin{itemize}\setlength\itemsep{0.8em}
    \item $\varphi(t)$ fits today's term structure.
    \item Two factors can represent level and slope-like movements.
    \item More parameters improve flexibility but complicate identification and validation.
  \end{itemize}
  \takeaway{Additional factors should answer a measured model deficiency, not merely add sophistication.}
  \src{QuantLib G2 model; portfolio.py.}
\end{frame}

\lesson{Calibration method becomes user-visible}
{The trade stores relative-price-error or price-error selection.}
{QuantLib helper enums are displayed rather than hidden behind an opaque label.}
{The model view reports fixed and calibrated parameters separately.}
{A user can explain why two calibrations differ instead of seeing only two NPVs.}
{portfolio.py calibration method normalization; trade form.}

\lesson{The model page documents the object graph}
{It names instruments, factories, signatures, purposes, dependencies, fields, and enumerations.}
{Schedules and model parameters become inspectable educational objects.}
{Research references are linked to the actual repository implementations.}
{Documentation is generated from the same normalized state used by pricing.}
{tgir\_quantlib\_tools/quantlib\_model.py; templates/quantlib\_model.html.}

\lesson{Workbook alignment adds an external benchmark}
{A second Bermudan trade is seeded with the spreadsheet reference economics.}
{The test expects a mark near USD 1.414 million within a USD 20,000 tolerance.}
{The schedule, surface interpolation, model, and tree must work together to meet the benchmark.}
{A benchmark test is stronger than checking only that NPV is positive.}
{tests/test\_sofr\_curve.py workbook reference test.}

\lesson{Exercise probabilities remain a known tree limitation}
{The UI lists the call schedule and calibration sources.}
{QuantLib's tree swaption engine returns value but does not expose the exercise probability profile used by this app.}
{The repository says ``not exposed'' rather than inventing statistics.}
{The later XCCY LSM engine records exercise counts explicitly because it owns the policy.}
{portfolio.py exercise schedule rows; README.md.}

\lesson{Sensitivity maps reveal calibration dependency}
{The European swaption responds to one selected matrix pillar.}
{The Bermudan responds to the diagonal and any interpolation source points.}
{Mean-reversion bumps rebuild and recalibrate the model.}
{Tests assert the exact nonzero volatility labels, turning data lineage into executable behavior.}
{tests/test\_sofr\_curve.py sensitivity tests.}

\lesson{The calibrated model still has deliberate simplifications}
{A single constant volatility parameter compresses an entire diagonal strip.}
{Mean reversion is fixed rather than statistically identified.}
{The one-factor model cannot reproduce all swaption-surface shapes simultaneously.}
{The workbench surfaces residuals so users can see the cost of parsimony.}
{portfolio.py calibration diagnostics.}

\lesson{Single-currency Bermudans are a tree problem here}
{One curve economy and one short-rate state keep lattice dimensionality manageable.}
{QuantLib already has the instrument, model, calibration helpers, and tree engine.}
{This native path is retained unchanged when cross-currency work begins.}
{A new XCCY engine should not destabilize a validated benchmark path.}
{XCCY pricing specification sections 1 and 9.}

\chronology{5 Apr 2026}{On-demand risk follows calibrated pricing}
{Commit \texttt{aa50775} adds roughly 688 lines for trade risk analytics and richer trade pages.}
{Curve, volatility, and model-parameter bumps are computed when requested.}
{The dashboard stays responsive by avoiding unconditional scenario expansion.}
{Risk is treated as a controlled workflow around the same pricing functions.}
{Commit aa50775.}

\lesson{Recalibration policy defines a Greek}
{A curve or volatility bump may change both direct valuation inputs and calibrated model parameters.}
{Holding parameters fixed measures one risk; full recalibration measures another.}
{The repository's diagnostics state which convention is used.}
{Comparisons are meaningful only when bump and calibration policies agree.}
{portfolio.py risk functions; xccy\_pricer\_diagnostics.py.}

\lesson{The Bermudan phase prepares the XCCY extension}
{It establishes schedule semantics, diagonal calibration, residual reporting, and explicit exercise dates.}
{It demonstrates native QuantLib backward induction in a low-dimensional setting.}
{It also exposes what the native engine does not report or generalize.}
{The later hybrid model keeps the economics and replaces only the state evolution and stopping engine.}
{Git chronology and XCCY specification.}

\lesson{A calibrated tree is a benchmark, not a universal engine}
{Its strengths are speed, deterministic backward induction, and native QuantLib integration.}
{Its boundary is the one-economy swaption instrument and low-dimensional short-rate state.}
{Cross-currency optionality adds a second rate factor, FX, quanto drift, and two-currency cash flows.}
{Recognizing that boundary prevents forcing the wrong abstraction onto the next product.}
{QuantLib source assessment; repository model architecture.}
